\documentclass[11pt,a4paper]{article}
\usepackage[utf8]{inputenc}
\usepackage{cmap}
\usepackage[T1]{fontenc}
\usepackage{lmodern}
\usepackage[margin=2.4cm]{geometry}
\usepackage{array,longtable,booktabs}
\usepackage[hidelinks]{hyperref}
\usepackage{amsmath}
\setlength{\parindent}{0pt}
\setlength{\parskip}{6pt}
\renewcommand{\contentsname}{Índice}
\emergencystretch=2em
\begin{document}
\begin{titlepage}
\centering
\vspace*{3cm}
{\Large Trabajo Práctico 3 de Compiladores\par}
\vspace{1cm}
{\LARGE\bfseries Analizador y Evaluador de Expresiones Trigonométricas\\[0.3cm] \textemdash{} Trabajo Práctico 3\par}
\vspace{2cm}
{\large Informe técnico\par}
\vspace{2cm}
\begin{tabular}{ll}
Estudiantes: & Juan rodriguez \\
             & Manuel Beleño \\
             & Gabriela Garcia \\
             & Josue fernandez \\[0.5cm]
Asignatura: & Compiladores \\
Profesor: & Esmeide Leal \\
Institución: & Universidad del Magdalena \\
Fecha: & 21 de noviembre de 2026 \\
\end{tabular}
\vfill
\end{titlepage}
\tableofcontents
\newpage
\section{Introducción}
El trabajo aborda el reconocimiento y evaluación de expresiones que combinan números reales, variables simbólicas, operaciones aritméticas y funciones trigonométricas. La cadena de procesamiento separa la identificación de símbolos, la validación de su estructura y el cálculo de su significado numérico. Esta separación permite detectar errores antes de solicitar valores y facilita la comprobación de cada componente.

La implementación está escrita en C99 y contiene solo los componentes del Trabajo Práctico 3: lexer, parser descendente, AST, evaluador, consola y verificador LL(1) de la gramática matemática.
\section{Objetivos}
\textbf{Objetivo general.} Diseñar e implementar un reconocedor LL(1) y evaluador de expresiones aritméticas y trigonométricas con resultados \texttt{double}.

\textbf{Objetivos específicos.} Definir una gramática y comprobar sus conjuntos de selección; reconocer expresiones mediante descenso recursivo; construir un AST; solicitar una vez el valor de cada variable; evaluar operaciones y funciones en radianes; detectar errores y validar los resultados mediante casos reproducibles.
\section{Diseño de la gramática}
La gramática final, en EBNF, usa \texttt{<Inicio>} como símbolo inicial. Las llaves representan repetición opcional y los corchetes una aparición opcional:
\begin{verbatim}
<Inicio>   -> <Expr>
<Expr>     -> <Termino> { ( "+" | "-" ) <Termino> }
<Termino>  -> <Unario> { ( "*" | "/" ) <Unario> }
<Unario>   -> "-" <Unario> | <Potencia>
<Potencia> -> <Primario> [ "^" <Unario> ]
<Primario> -> NUMBER | ID | "(" <Expr> ")"
              | <Funcion> "(" <Expr> ")"
<Funcion>  -> "sin" | "sen" | "cos" | "tan"
\end{verbatim}
Los no terminales se escriben entre ángulos. Los terminales son \texttt{NUMBER}, \texttt{ID}, los operadores, paréntesis y los cuatro nombres de función. \texttt{ID} designa una variable simbólica; los nombres de función no se clasifican como identificadores. El programa exige fin de entrada después de la expresión.

La precedencia aumenta desde suma y resta, multiplicación y división, negación unaria, hasta potencia y primarios. Las primeras operaciones se agrupan por la izquierda. La potencia se agrupa por la derecha: \verb|2^3^2| equivale a \verb|2^(3^2)|. La potencia precede a la negación: \verb|-2^2| equivale a \verb|-(2^2)|; también se admite \verb|2^-2|. Los paréntesis alteran el agrupamiento. \texttt{sen} es un alias de \texttt{sin}; todas las funciones reciben una expresión entre paréntesis y operan en radianes.
\section{Verificación LL(1)}
El verificador \texttt{verificar\_ll1} expande opciones y repeticiones a producciones auxiliares, calcula anulabilidad, PRIMERO y SIGUIENTE por punto fijo y compara los conjuntos de selección de las alternativas. La orden \verb|./verificar_ll1 gramatica/expresiones_matematicas.txt --sets| produjo \texttt{ninguno}, es decir, cero conflictos detectados. Se presupone la clasificación léxica disjunta de las funciones e \texttt{ID}.

El programa devuelve 0 sin conflictos, 1 ante conflictos o errores de lectura, 2 para una gramática vacía o mal formada y 3 ante un fallo de reserva o tamaño no representable. Los fallos de memoria interrumpen el análisis y liberan las estructuras parciales. La suite \texttt{test\_ll1} comprueba mediante la API C la gramática matemática, las gramáticas inválidas y la detección de conflictos.

Los conjuntos principales obtenidos son los siguientes. Se escribe \texttt{\$} para fin de entrada y F para \{sin, sen, cos, tan\}.
\begin{center}
\begin{tabular}{lll}\toprule
No terminal & PRIMERO & SIGUIENTE \\\midrule
Inicio & \{ (, -, F, ID, NUMBER \} & \{ \$ \} \\
Expr & \{ (, -, F, ID, NUMBER \} & \{ ), \$ \} \\
Termino & \{ (, -, F, ID, NUMBER \} & \{ ), +, -, \$ \} \\
Unario & \{ (, -, F, ID, NUMBER \} & \{ ), *, /, +, -, \$ \} \\
Potencia & \{ (, F, ID, NUMBER \} & \{ ), *, /, +, -, \$ \} \\
Primario & \{ (, F, ID, NUMBER \} & \{ ), *, /, +, -, \verb|^|, \$ \} \\
Funcion & F & \{ ( \} \\\bottomrule
\end{tabular}
\end{center}
Las alternativas de primario tienen comienzos distintos, y los conjuntos de selección de las reglas auxiliares son disjuntos. La tabla predictiva completa no es necesaria para el reconocedor recursivo utilizado.
\section{Análisis léxico}
\texttt{lexer\_init} configura el único modo matemático. El lector reconoce operadores, paréntesis, números, identificadores y las cuatro funciones, con línea y columna. Los enteros y decimales de una expresión siguen el patrón \verb|[0-9]+(\.[0-9]+)?|; no se aceptan \texttt{.5}, \texttt{5.} ni exponentes en literales de expresión. Los identificadores admiten letras ASCII, dígitos tras el primer carácter y guion bajo. El reconocimiento de máxima coincidencia distingue \texttt{sin} de \texttt{sinus}. La secuencia \verb|--| son dos signos negativos, no un comentario. Un carácter inválido o un decimal mal formado produce un token de error.
\section{Análisis sintáctico}
\texttt{parser\_expr.c} utiliza un token de preanálisis. Las funciones \texttt{expr}, \texttt{termino}, \texttt{unario}, \texttt{potencia}, \texttt{primario} y \texttt{funcion} corresponden a las reglas de la gramática. Los bucles de suma y producto construyen nodos sucesivos por la izquierda; la recursión del exponente construye el subárbol derecho. El punto de entrada exige EOF. Los errores informan posición, token encontrado y elemento esperado cuando corresponde. Se rechazan, entre otros, \texttt{sin()}, \texttt{2+} y texto sobrante.
\section{Árbol de sintaxis abstracta}
El tipo \texttt{ASTExpr} modela número e identificador con su texto; los paréntesis y la negación contienen un hijo; el binario contiene operador e hijos; la llamada contiene nombre y argumentos. Por ejemplo, \verb|2+3*4| produce una suma cuyo hijo derecho es una multiplicación. Los operadores se almacenan como cadenas tales como \texttt{TOK-PLUS} y \texttt{TOK-EXP}. La arena y los nodos pertenecen al parser y son válidos hasta \texttt{parser\_expr\_destroy}.
\section{Evaluación semántica}
El evaluador recorre el AST sin modificarlo: convierte números con \texttt{strtod}, consulta variables mediante un callback, evalúa hijos recursivamente y aplica el operador o función correspondiente. Usa \texttt{double}, \texttt{pow}, \texttt{sin}, \texttt{cos} y \texttt{tan}; verifica conversión, rango, dominio, división por cero y valores no finitos. El resultado distingue explícitamente éxito y error.

El enunciado propone calcular durante el descenso sintáctico. El diseño acordado realiza la evaluación semántica \emph{después} de validar toda la expresión y construir el AST. Esta adaptación permite solicitar variables únicamente tras confirmar la sintaxis y reutilizar el árbol sin duplicar el parser. La evaluación sigue la estructura construida por el descenso recursivo.

El descenso sintáctico tiene un límite de 128 niveles de \texttt{<Unario>}; la recolección de variables y la evaluación tienen un límite independiente de 128 nodos de profundidad en el AST. Al superar cualquiera se informa un error controlado.
\section{Interfaz de consola}
\texttt{calculadora\_expr} fija \texttt{LC\_NUMERIC} en \texttt{C}, lee una expresión por línea y la analiza. Si es válida, recorre el AST de izquierda a derecha, registra cada variable distinta y solicita su valor una sola vez. Los valores inválidos permiten reintento; EOF termina la ejecución. Después evalúa y muestra el resultado. Se liberan la lista de variables, la arena del parser y la línea. Los valores de variables se convierten con \texttt{strtod} y admiten signo y notación científica, a diferencia de los literales de la gramática.

La compilación usa solo C99 y la biblioteca estándar. En Linux se emplea \texttt{make}; en Windows se emplea \texttt{build.bat} con MinGW-w64/GCC desde \texttt{cmd}. El script no requiere interpretes, MSYS shell ni enlazadores GNU. Los mensajes de ejecución son ASCII para evitar dependencias de la pagina de codigos de consola.

El modo \texttt{--tokens archivo.txt} ejecuta únicamente el analizador léxico y escribe un token por línea con su tipo y lexema. No invoca al parser ni al evaluador, por lo que sirve para inspeccionar la tokenización de forma independiente. Devuelve 0 por éxito, 1 ante un error de archivo o memoria y 2 ante un error léxico, que se informa con línea y columna.
\section{Casos de prueba y resultados}
Se ejecutaron los siete ejemplos del enunciado con el programa de consola y se verificaron mediante la suite C \texttt{test\_consola\_expr}, con tolerancia absoluta y relativa de $10^{-12}$. Los valores esperados son constantes fijas precalculadas dentro de los tests; no se calculan con \texttt{sin}, \texttt{cos} ni \texttt{tan} durante la prueba. La identidad trigonométrica se compara con la constante $1$. Los resultados siguientes son los observados, redondeados a 15 cifras; el documento \texttt{casos\_prueba\_validacion.md} contiene los valores con mayor precisión y más pruebas.
\begin{longtable}{p{5.7cm}p{2.5cm}p{3.8cm}l}\toprule
Expresión & Variables & Resultado & Estado \\\midrule\endhead
\verb|sin(3.14/2)+cos(0)*2^2| & -- & 4.99999968293183 & Pasa \\
\verb|sin(x)+cos(y)| & x=0, y=0 & 1 & Pasa \\
\verb|sin(cos(x))+tan(y)/cos(x)| & x=0, y=0 & 0.841470984807897 & Pasa \\
\verb|cos(x)^2+sin(x)^2| & x=1 & 1 & Pasa \\
\verb|sin(x)| & x=1 & 0.841470984807897 & Pasa \\
\verb|cos(x+y)| & x=1, y=2 & -0.989992496600445 & Pasa \\
\verb|tan(x^2+sin(x))| & x=1 & -3.60380353200881 & Pasa \\
\verb|2+3*4| & -- & 14 & Pasa \\
\verb|2^3^2| & -- & 512 & Pasa \\
\verb|-2^2| & -- & -4 & Pasa \\
\verb|2^-2| & -- & 0.25 & Pasa \\
\verb|x+x| & x=3 & 6 & Pasa \\\bottomrule
\end{longtable}
El primer ejemplo es aproximadamente 4.99999968293183, no exactamente 5, porque 3.14 aproxima $\pi$. La identidad $\cos^2(x)+\sin^2(x)\approx 1$ pasó para $x=-7,-1,0,0.5,3.14,100$.

\begin{center}\begin{tabular}{lll}\toprule
Entrada & Resultado observado & Estado \\\midrule
\texttt{sin()}, \texttt{2+}, \texttt{cos(x} & Error sintáctico, código 2 & Pasa \\
\texttt{sin(1,2)} & Error léxico, código 2 & Pasa \\
\texttt{1/0}, \verb|0^-1| & Error matemático, código 4 & Pasa \\
\verb|(-1)^0.5|, \verb|10^1000| & Dominio o rango, código 4 & Pasa \\\bottomrule\end{tabular}\end{center}
Se ejecutaron \texttt{make clean}, \texttt{make all} y \texttt{make test} con \texttt{-Werror}. Pasaron las siete suites C: parser, evaluador, consola, tokens, API LL(1), CLI de tokens y CLI del verificador. \texttt{test\_tokens} comprueba el volcado léxico y los errores de archivo; \texttt{test\_tokens\_cli} ejecuta \texttt{calculadora\_expr --tokens}. \texttt{test\_ll1} lee el archivo real de la gramática; \texttt{test\_ll1\_cli} ejecuta \texttt{verificar\_ll1} con conjuntos, un conflicto, tres gramáticas inválidas y un archivo inexistente. La verificación de Windows se realiza con \texttt{build.bat all} y \texttt{build.bat test}; el entorno de este informe no dispone de MinGW para ejecutarla. No se midió cobertura.
\section{Conclusiones}
Se obtuvo un reconocedor LL(1), un AST y un evaluador numérico integrados en una consola. La separación de fases permitió identificar errores, verificar la precedencia mediante la estructura del árbol y mantener un único conjunto de pruebas C portable.

Las limitaciones incluyen una expresión por ejecución, multiplicación explícita, literales de expresión sin notación científica y los límites explícitos de profundidad. Los errores \texttt{ERANGE} se rechazan también en subdesbordamiento; no se aplican tolerancias arbitrarias a singularidades de la tangente.
\section{Referencias}
\begin{enumerate}
\item Enunciado suministrado: \texttt{docs/compiladores\_trabajo\_pr\_ctico\_3.md}.
\item Implementación y documentación del repositorio: \path{gramatica/expresiones_matematicas.txt}, \path{src/}, \path{tests/} y \path{docs/gramatica_expresiones_matematicas.md}.
\item Resultados experimentales reproducibles: \path{tests/} y \path{docs/casos_prueba_validacion.md}.
\end{enumerate}
\end{document}
